\documentclass[a4paper,12pt,oneside]{article}
\makeatletter
\def\input@path{{../../}}
\makeatother
\usepackage{discrete-practice}

\begin{document}

\PracticeHeader{2}{1}{Графы}

\Task{Задание 1}
Может ли быть простым граф, содержащий 4 вершины и 8 ребер?
Может ли граф с одним ребром быть псевдографом?
Может ли граф быть псевдографом, если в нем нет кратных ребер?

\Task{Задание 2}
Выяснить, существуют ли графы с набором степеней:
\[
  1)\;0,2,2,3,3\qquad 2)\;2,2,2,3,3\qquad
  3)\;2,2,3,3,3\qquad 4)\;0,1,2,3,4 .
\]
Построить, если существуют, если нет -- объясните почему.

2) Степени вершин графа: $1,2,3,4,5,6,7,x$. Чему может быть равен $x$?

\Task{Задание 3}
Определить число ребер в каждом из графов $K_n$, $K_{p,q}$.

\Task{Задание 4}
В полном графе $18$ вершин. Сколько в нем ребер, инцидентных одной вершине?

\Task{Задание 5}
Полный граф имеет $105$ ребер. Найдите число его вершин.

\Task{Задание 6}
При каких $n$ существуют простые графы с $n$ вершинами, каждая из которых
имеет степень 3? степень 4?

\clearpage

\Task{Задание 7}
Графы, изображенные на рисунке, разбить на классы попарно изоморфных графов.
\begin{center}
  \includegraphics[width=.76\linewidth]{Practice_1_1.png}
\end{center}

\Task{Задание 7.1}
Постройте 2 неизоморфных графа с $|V|=6$ и $|E|=8$.

\Task{Задание 8}
\begin{center}
  \includegraphics[width=.58\linewidth]{Practice_1_2.png}
\end{center}
\begin{numlist}
  \item Укажите вершины, инцидентные ребру $a$.
  \item Укажите номера вершин, содержащих петли.
\end{numlist}

\end{document}
